\ifdefined\gkver\else\def\gkver{student}\fi
\documentclass[\gkver]{gaokaozhenti}
\begin{document}

\chapter{1997年上海}

\begin{enumerate}
\item
%% number: 1
%% paperName: 1997年上海高考第 1 题 4 分
%% typeId: 1
%% type: 单选题
%% chapter: 机械振动
%% point: 弹簧振子
%% method: 无
%% score: 4
%% degree: 650
%% duplicateId: 0
%% body:
弹簧振子在光滑水平面上作简谐振动。在振子向平衡位置运动的过程中\xzanswer{D}
\fourchoices[answer=D]
{振子所受的回复力逐渐增大}
{振子的位移逐渐增大}
{振子的速度逐渐减小}
{振子的加速度逐渐减小}

%% answer: D
%% memo:
\memoanswer{
本题原卷及材料中未提供详解，待补充。
}

\item
%% number: 2
%% paperName: 1997年上海高考第 2 题 4 分
%% typeId: 1
%% type: 单选题
%% chapter: 气体性质
%% point: 理想气体状态方程
%% method: 无
%% score: 4
%% degree: 650
%% duplicateId: 0
%% body:
对一定质量的理想气体，下列说法正确的是\xzanswer{A}
\fourchoices[answer=A]
{压强增大，体积增大，分子的平均动能一定增大}
{压强减小，体积减小，分子的平均动能一定增大}
{压强减小，体积增大，分子的平均动能一定增大}
{压强增大，体积减小，分子的平均动能一定增大}

%% answer: A
%% memo:
\memoanswer{
本题原卷及材料中未提供详解，待补充。
}

\item
%% number: 3
%% paperName: 1997年上海高考第 3 题 4 分
%% typeId: 1
%% type: 单选题
%% chapter: 原子和原子核
%% point: 质能方程
%% method: 无
%% score: 4
%% degree: 650
%% duplicateId: 0
%% body:
氘核（\ce{^{2}_{1}H}）和氚核（\ce{^{3}_{1}H}）聚合成氦核（\ce{^{4}_{2}He}）的核反应方程如下：

\ce{^{2}_{1}H + ^{3}_{1}H -> ^{4}_{2}He + ^{1}_{0}n}

设氘核质量为 $m_{1}$，氚核质量为 $m_{2}$，氦核质量为 $m_{3}$，中子质量为 $m_{4}$，则反应过程中释放的能量为\xzanswer{C}
\fourchoices[answer=C]
{（$m_{1} + m_{2} - m_{3}$）$c^{2}$}
{（$m_{1} + m_{2} - m_{4}$）$c^{2}$}
{（$m_{1} + m_{2} - m_{3} - m_{4}$）$c^{2}$}
{（$m_{3} + m_{4} - m_{1} - m_{2}$）$c^{2}$}

%% answer: C
%% memo:
\memoanswer{
本题原卷及材料中未提供详解，待补充。
}

\item
%% number: 4
%% paperName: 1997年上海高考第 4 题 4 分
%% typeId: 1
%% type: 单选题
%% chapter: 电磁感应
%% point: 电磁感应中图象
%% method: 无
%% score: 4
%% degree: 650
%% duplicateId: 0
%% body:
一磁棒自远处匀速沿一圆形线圈的轴线运动，并穿过线圈向远处而去，如右图所示。则下列四图中，较正确反映线圈中电流 $i$ 与时间 $t$ 关系的是（线圈中电流以图示箭头为正方向）\xzanswer{B}
\onepicture[w=5cm]{\includesvg{figs/04a}}
\fourchoices[answer=B, ispicture=true, h=2.2cm]
{figs/04b.pdf}
{figs/04c.pdf}
{figs/04d.pdf}
{figs/04e.pdf}

%% answer: B
%% memo:
\memoanswer{
本题原卷及材料中未提供详解，待补充。
}

\item
%% number: 5
%% paperName: 1997年上海高考第 5 题 4 分
%% typeId: 1
%% type: 单选题
%% chapter: 曲线运动
%% point: 竖直平面圆周运动
%% method: 无
%% score: 4
%% degree: 450
%% duplicateId: 0
%% body:
质量为 $m$ 的小球被系在轻绳一端，在竖直平面内作半径为 $R$ 的圆周运动，运动过程中小球受到空气阻力的作用。设某一时刻小球通过轨道的最低点，此时绳子的张力为 $7mg$，此后小球继续作圆周运动，经过半个圆周恰能通过最高点，则在此过程中小球克服空气阻力所做的功为\xzanswer{C}
\fourchoices[answer=C]
{$\frac{1}{4}mgR$}
{$\frac{1}{3}mgR$}
{$\frac{1}{2}mgR$}
{$mgR$}

%% answer: C
%% memo:
\memoanswer{
本题原卷及材料中未提供详解，待补充。
}

\item
%% number: 6
%% paperName: 1997年上海高考第 6 题 4 分
%% typeId: 1
%% type: 单选题
%% chapter: 力物体的平衡
%% point: 力矩平衡
%% method: 无
%% score: 4
%% degree: 450
%% duplicateId: 0
%% body:
如图所示是一种手控制动器。a 是一个转动着的轮子，b 是摩擦制动片，c 是杠杆，O 是其固定转动轴。手在 A 点施加一个作用力 $F$ 时，b 将压紧轮子，使轮子制动。若使轮子制动所需的力矩是一定的，则下列说法正确的是\xzanswer{A}
\onepicture[h=4.5cm]{\includesvg{figs/06}}
\fourchoices[answer=A]
{轮 a 逆时针转动时，所需的力 $F$ 较小}
{轮 a 顺时针转动时，所需的力 $F$ 较小}
{无论 a 逆时针还是顺时针转动，所需的力 $F$ 相同}
{无法比较 $F$ 的大小}

%% answer: A
%% memo:
\memoanswer{
本题原卷及材料中未提供详解，待补充。
}

\item
%% number: 7
%% paperName: 1997年上海高考第 7 题 5 分
%% typeId: 2
%% type: 多选题
%% chapter: 机械波
%% point: 波的图象
%% method: 无
%% score: 5
%% degree: 650
%% duplicateId: 0
%% body:
一列横波沿水平方向传播，某一时刻的波形如图所示，则图中 a、b、c、d 四点在此时刻具有相同运动方向的是\xzanswer{BC}
\onepicture[w=6cm]{\includesvg{figs/07}}
\fourchoices[answer=BC]
{a 和 c}
{a 和 d}
{b 和 c}
{b 和 d}

%% answer: BC
%% memo:
\memoanswer{
本题原卷及材料中未提供详解，待补充。
}

\item
%% number: 8
%% paperName: 1997年上海高考第 8 题 5 分
%% typeId: 2
%% type: 多选题
%% chapter: 运动和力
%% point: 动量守恒定律
%% method: 无
%% score: 5
%% degree: 650
%% duplicateId: 0
%% body:
在光滑水平面上，两球沿球心连线以相等速率相向而行，并发生碰撞，下列现象可能的是\xzanswer{AD}
\fourchoices[answer=AD]
{若两球质量相同，碰后以某一相等速率互相分开}
{若两球质量相同，碰后以某一相等速率同向而行}
{若两球质量不同，碰后以某一相等速率互相分开}
{若两球质量不同，碰后以某一相等速率同向而行}

%% answer: AD
%% memo:
\memoanswer{
本题原卷及材料中未提供详解，待补充。
}

\item
%% number: 9
%% paperName: 1997年上海高考第 9 题 5 分
%% typeId: 2
%% type: 多选题
%% chapter: 磁场
%% point: 左手定则
%% method: 无
%% score: 5
%% degree: 650
%% duplicateId: 0
%% body:
如图所示，一金属直杆 MN 两端接有导线，悬挂于线圈上方，MN 与线圈轴线均处于竖直平面内，为使 MN 垂直纸面向外运动，可以\xzanswer{ABD}
\onepicture[w=5cm]{\includesvg{figs/09}}
\fourchoices[answer=ABD]
{将 a、c 端接在电源正极，b、d 端接在电源负极}
{将 b、d 端接在电源正极，a、c 端接在电源负极}
{将 a、d 端接在电源正极，b、c 端接在电源负极}
{将 a、c 端接在交流电源的一端，b、d 接在交流电源的另一端}

%% answer: ABD
%% memo:
\memoanswer{
本题原卷及材料中未提供详解，待补充。
}

\item
%% number: 10
%% paperName: 1997年上海高考第 10 题 5 分
%% typeId: 2
%% type: 多选题
%% chapter: 运动和力
%% point: 变加速直线运动
%% method: 无
%% score: 5
%% degree: 450
%% duplicateId: 0
%% body:
某人身系弹性绳自高空 P 点自由下落，图中 a 点是弹性绳的原长位置，c 是人所到达的最低点，b 是人静止地悬吊着时的平衡位置，不计空气阻力，则下列说法中正确的是\xzanswer{BC}
\onepicture[h=5cm]{\includesvg{figs/10}}
\fourchoices[answer=BC]
{从 P 至 c 过程中重力的冲量大于弹性绳弹力的冲量}
{从 P 至 c 过程中重力所做功等于人克服弹力所做的功}
{从 P 至 b 过程中人的速度不断增大}
{从 a 至 c 过程中加速度方向保持不变}

%% answer: BC
%% memo:
\memoanswer{
本题原卷及材料中未提供详解，待补充。
}

\item
%% number: 11
%% paperName: 1997年上海高考第 11 题 5 分
%% typeId: 2
%% type: 多选题
%% chapter: 电场
%% point: 带电粒子的直线运动
%% method: 无
%% score: 5
%% degree: 450
%% duplicateId: 0
%% body:
如图所示，A、B 为平行金属板，两板相距为 $d$，分别与电源两极相连，两板的中央各有一小孔 M 和 N。今有一带电质点，自 A 板上方相距为 $d$ 的 P 点由静止自由下落（P、M、N 在同一竖直线上），空气阻力忽略不计，到达 N 孔时速度恰好为零，然后沿原路返回。若保持两极板间的电压不变，则\xzanswer{ACD}
\onepicture[w=6cm]{\includesvg{figs/11}}
\fourchoices[answer=ACD]
{把 A 板向上平移一小段距离，质点自 P 点自由下落后仍能返回}
{把 A 板向下平移一小段距离，质点自 P 自由下落后将穿过 N 孔继续下落}
{把 B 板向上平移一小段距离，质点自 P 点自由下落后仍能返回}
{把 B 板向下平移一小段距离，质点自 P 点自由下落后将穿过 N 孔继续下落}

%% answer: ACD
%% memo:
\memoanswer{
本题原卷及材料中未提供详解，待补充。
}

\item
%% number: 12
%% paperName: 1997年上海高考第 12 题 4 分
%% typeId: 3
%% type: 填空题
%% chapter: 原子和原子核
%% point: 原子核式结构
%% method: 无
%% score: 4
%% degree: 650
%% duplicateId: 0
%% body:
（A 组）关于原子结构，卢瑟福根据 \tkanswer[2.6cm]{$\alpha$粒子散射} 实验提出了原子核式结构模型。本世纪初，玻尔结合普朗克量子学说提出的原子模型，成功地解释了 \tkanswer[1.2cm]{氢} 原子的光谱。

%% answer: 
\jdanswer{
$\alpha$粒子散射，氢
}
%% memo:
\memoanswer{
本题原卷及材料中未提供详解，待补充。
}

\item
%% number: 12
%% paperName: 1997年上海高考第 12 题 4 分
%% typeId: 3
%% type: 填空题
%% chapter: 光学
%% point: 光电效应
%% method: 无
%% score: 4
%% degree: 650
%% duplicateId: 0
%% body:
（B 组）关于光的本性，早期有牛顿的微粒说和惠更斯的 \tkanswer[1.6cm]{波动} 说，后来又有麦克斯韦的电磁说。本世纪初，为解释 \tkanswer[1.8cm]{光电效应} 现象，爱因斯坦提出了光子说。

%% answer: 
\jdanswer{
波动，光电效应
}
%% memo:
\memoanswer{
本题原卷及材料中未提供详解，待补充。
}

\item
%% number: 13
%% paperName: 1997年上海高考第 13 题 4 分
%% typeId: 3
%% type: 填空题
%% chapter: 曲线运动
%% point: 平抛运动
%% method: 无
%% score: 4
%% degree: 650
%% duplicateId: 0
%% body:
（A 组）在一次“飞车过黄河”的表演中，汽车在空中飞经最高点后在对岸着地。已知汽车从最高点至着地点经历时间约 $0.8\Us$，两点间的水平距离约为 $30\Um$，忽略空气阻力，则最高点与着地点间的高度差约为 \tkanswer[1.4cm]{3.2} $\Um$。

%% answer: 
\jdanswer{
3.2
}
%% memo:
\memoanswer{
本题原卷及材料中未提供详解，待补充。
}

\item
%% number: 13
%% paperName: 1997年上海高考第 13 题 4 分
%% typeId: 3
%% type: 填空题
%% chapter: 曲线运动
%% point: 平抛运动
%% method: 无
%% score: 4
%% degree: 650
%% duplicateId: 0
%% body:
（B 组）在一次“飞车过黄河”的表演中，汽车在空中飞经最高点后在对岸着地。已知汽车从最高点至着地点经历时间约 $0.8\Us$，两点间的水平距离约为 $30\Um$，忽略空气阻力，则汽车在最高点时的速度约为 \tkanswer[1.6cm]{37.5} $\Ums$。

%% answer: 
\jdanswer{
37.5
}
%% memo:
\memoanswer{
本题原卷及材料中未提供详解，待补充。
}

\item
%% number: 14
%% paperName: 1997年上海高考第 14 题 4 分
%% typeId: 3
%% type: 填空题
%% chapter: 交流电
%% point: 交流电
%% method: 无
%% score: 4
%% degree: 650
%% duplicateId: 0
%% body:
（A 组）一交流电压的瞬时值表达式为 $u = 15\sin 100\pi t\ (\UV)$，将该交流电压加在一电阻上时，产生的电功率为 $25\UW$。那么，这个电阻的阻值是 \tkanswer[1.4cm]{4.5} $\UO$。

%% answer: 
\jdanswer{
$4.5\UO$
}
%% memo:
\memoanswer{
本题原卷及材料中未提供详解，待补充。
}

\item
%% number: 14
%% paperName: 1997年上海高考第 14 题 4 分
%% typeId: 3
%% type: 填空题
%% chapter: 电磁感应
%% point: 远距离输电
%% method: 无
%% score: 4
%% degree: 650
%% duplicateId: 0
%% body:
（B 组）水电站给远处山村送电的输出功率是 $100\UkW$，用 $2000\UV$ 电压输电，线路上损失的功率是 $2.5\times10^{4}\UW$，如果改用 $20000\UV$ 高压输电，线路上损失的功率是 \tkanswer[1.6cm]{250} $\UW$。

%% answer: 
\jdanswer{
$250\UW$
}
%% memo:
\memoanswer{
本题原卷及材料中未提供详解，待补充。
}

\item
%% number: 15
%% paperName: 1997年上海高考第 15 题 4 分
%% typeId: 3
%% type: 填空题
%% chapter: 交流电
%% point: LC振荡回路
%% method: 无
%% score: 4
%% degree: 650
%% duplicateId: 0
%% body:
已知 $LC$ 振荡电路中电感为 $L_{0}$，为利用此振荡电路发射波长为 $\lambda$ 的电磁波，振荡电路的电容应为 \tkanswer[3.2cm]{$\frac{\lambda^{2}}{4\pi^{2}c^{2}L_{0}}$}。

%% answer: 
\jdanswer{
$\frac{\lambda^{2}}{4\pi^{2}c^{2}L_{0}}$
}
%% memo:
\memoanswer{
本题原卷及材料中未提供详解，待补充。
}

\item
%% number: 16
%% paperName: 1997年上海高考第 16 题 4 分
%% typeId: 3
%% type: 填空题
%% chapter: 电磁感应
%% point: 电磁感应的电量
%% method: 无
%% score: 4
%% degree: 450
%% duplicateId: 0
%% body:
如图所示，空间存在垂直于纸面的均匀磁场，在半径为 $a$ 的圆形区域内、外，磁场方向相反，磁感应强度的大小均为 $B$。一半径为 $b$，电阻为 $R$ 的圆形导线环放置在纸面内，其圆心与圆形区域的中心重合。在内、外磁场同时由 $B$ 均匀地减小到零的过程中，通过导线截面的电量 $q =$ \tkanswer[3.2cm]{$\frac{B\pi(2a^{2}-b^{2})}{R}$（或 $\frac{B\pi(b^{2}-2a^{2})}{R}$）}。
\onepicture[align=right, w=4.5cm]{\includesvg{figs/16}}

%% answer: 
\jdanswer{
$\frac{B\pi(2a^{2}-b^{2})}{R}$（或 $\frac{B\pi(b^{2}-2a^{2})}{R}$）
}
%% memo:
\memoanswer{
本题原卷及材料中未提供详解，待补充。
}

\item
%% number: 17
%% paperName: 1997年上海高考第 17 题 4 分
%% typeId: 3
%% type: 填空题
%% chapter: 气体性质
%% point: 阿伏伽德罗常数
%% method: 无
%% score: 4
%% degree: 650
%% duplicateId: 0
%% body:
已知铜的密度为 $8.9\times10^{3}\Ukgmc$，铜的原子量为 64，质子和中子的质量均约为 $1.67\times10^{-27}\Ukg$，则铜块中平均每个铜原子所占的空间体积约为 \tkanswer[2.4cm]{$1.2\times10^{-29}$} $\Umc$。

%% answer: 
\jdanswer{
$1.2\times10^{-29}\Umc$
}
%% memo:
\memoanswer{
本题原卷及材料中未提供详解，待补充。
}

\item
%% number: 18
%% paperName: 1997年上海高考第 18 题 4 分
%% typeId: 4
%% type: 实验题
%% chapter: 光学
%% point: 光电效应
%% method: 无
%% score: 4
%% degree: 650
%% duplicateId: 0
%% body:
如图所示，一静电计与锌板相连，在 A 处用一紫外灯照射锌板，关灯后，指针保持一定偏角。
\begin{steps}
	\item
	现用一带负电的金属小球与锌板接触，则静电计指针偏角将 \tkanswer[1.4cm]{减小}（填“增大”，“减小”或“不变”）。

	\item
	使静电计指针回到零，再用相同强度的钠灯发出的黄光照射锌板，静电计指针无偏转。那么，若改用强度更大的红外线灯照射锌板，可观察到静电计指针 \tkanswer[1.4cm]{无}（填“有”或“无”）偏转。
\end{steps}
\onepicture[align=right, w=6cm]{\includesvg{figs/18}}

%% answer: 
\jdanswer{
\begin{enumerate}
	\item
	减小

	\item
	无

\end{enumerate}
}
%% memo:
\memoanswer{
本题原卷及材料中未提供详解，待补充。
}

\item
%% number: 19
%% paperName: 1997年上海高考第 19 题 6 分
%% typeId: 4
%% type: 实验题
%% chapter: 电路
%% point: 多用表的使用
%% method: 无
%% score: 6
%% degree: 650
%% duplicateId: 0
%% body:
用万用表测电阻。
\begin{steps}
	\item
	每次换档后，需重新 \tkanswer[1.4cm]{调零} 再进行测量。

	\item
	如果表的指针偏转过大，为使测量比较精确，应将选择开关拨至倍率较 \tkanswer[1.4cm]{小}（填大或小）的档位上。

	\item
	某次测电阻时，表的指针位置如图所示，则该电阻的测量值是 \tkanswer[1.4cm]{190} $\UO$。
\end{steps}
\onepicture[align=right, w=4.5cm]{\includesvg{figs/19}}

%% answer: 
\jdanswer{
\begin{enumerate}
	\item
	调零

	\item
	小

	\item
	$190\UO$

\end{enumerate}
}
%% memo:
\memoanswer{
本题原卷及材料中未提供详解，待补充。
}

\item
%% number: 20
%% paperName: 1997年上海高考第 20 题 6 分
%% typeId: 4
%% type: 实验题
%% chapter: 电路
%% point: 电路实验
%% method: 等效替代
%% score: 6
%% degree: 450
%% duplicateId: 0
%% body:
某同学用以下器材测量电阻：

①安培计、②电阻箱、③单刀双掷电键（这种电键在掷刀 a 倒向 b 时 ab 接通，倒向 c 时 ac 接通）、④待测电阻、⑤电源、⑥限流电阻，如图所示。实验方法是利用单刀双掷电键分别将电阻箱和待测电阻接入电路，用电阻箱替代待测电阻的方法来测定待测电阻的阻值。
\begin{enumerate}
	\item
	在图中完成电路的连接（其中有二条导线已连接好）。

	\item
	本实验中电键应先接通的是含有 \tkanswer[2.2cm]{待测电阻}（填“待测电阻”或“电阻箱”）的电路。用电键变换电路，调节电阻箱时，应使两次测量的 \tkanswer[2.8cm]{电流（或安培计读数）} 大小相同，这时待测电阻的值可以从 \tkanswer[2cm]{电阻箱} 上读出。
\end{enumerate}
\onepicture[align=right, w=7cm]{\includesvg{figs/20}}

%% answer: 
\jdanswer{
\begin{enumerate}
	\item
	如图
	\onepicture[w=7cm]{\includesvg{figs/20_an}}

	\item
	待测电阻，电流（或安培计读数），电阻箱

\end{enumerate}
}
%% memo:
\memoanswer{
本题原卷及材料中未提供详解，待补充。
}

\item
%% number: 21
%% paperName: 1997年上海高考第 21 题 6 分
%% typeId: 4
%% type: 实验题
%% chapter: 直线运动
%% point: 打点计时器公式
%% method: 无
%% score: 6
%% degree: 650
%% duplicateId: 0
%% body:
下图表示用打点计时器记录小车的运动情况，开始时小车在光滑水平玻璃板上运动。后来在薄布面上作匀减速运动，所打出的纸带如图所示（附有刻度尺），纸带上相邻两点对应的时间间隔为 $0.02\Us$。

从纸带上可以确定小车作匀减速运动的初速度是 \tkanswer[1.4cm]{0.9} $\Ums$，小车在布上运动的加速度大小是 \tkanswer[1.4cm]{5} $\Umsq$。
\onepicture[align=right, w=12cm]{\includesvg{figs/21}}

%% answer: 
\jdanswer{
$0.9\Ums$，$5\Umsq$
}
%% memo:
\memoanswer{
本题原卷及材料中未提供详解，待补充。
}

\item
%% number: 22
%% paperName: 1997年上海高考第 22 题 7 分
%% typeId: 4
%% type: 实验题
%% chapter: 运动和力
%% point: 牛顿第二定律
%% method: 无
%% score: 7
%% degree: 450
%% duplicateId: 0
%% body:
为测定木块与斜面之间的滑动摩擦系数，某同学让木块从斜面上端自静止起作匀加速下滑运动，他使用的实验器材仅限于①倾角固定的斜面（倾角未知）、②木块、③秒表、④米尺。

实验中应记录的数据是 \tkanswer[2.4cm]{$h$、$d$、$L$、$t$}，计算摩擦系数的公式是 $\mu =$ \tkanswer[3.6cm]{$\frac{h}{d}-\frac{2L^{2}}{gt^{2}d}$}。为了减小测量的误差，可采用的办法是 \tkanswer[3.6cm]{多次测量取平均值的方法}。
\onepicture[align=right, w=5.5cm]{\includesvg{figs/22}}

%% answer: 
\jdanswer{
$h$、$d$、$L$、$t$；$\frac{h}{d}-\frac{2L^{2}}{gt^{2}d}$；多次测量取平均值的方法
}
%% memo:
\memoanswer{
本题原卷及材料中未提供详解，待补充。
}

\item
%% number: 23
%% paperName: 1997年上海高考第 23 题 10 分
%% typeId: 6
%% type: 计算题
%% chapter: 光学
%% point: 光的折射
%% method: 无
%% score: 10
%% degree: 650
%% duplicateId: 0
%% body:
单色细光束射到折射率 $n = \sqrt{2}$ 的透明球表面，光束在过球心的平面内，入射角 $i = 45^{\circ}$。研究经折射进入球内后又经内表面反射一次，再经球面折射后射出的光线，如图所示（图上已画出入射光和出射光）。
\begin{enumerate}
	\item
	在图上大致画出光线在球内的路径和方向；
	\item
	求入射光与出射光之间的夹角 $\alpha$；
	\item
	如果入射的是一束白光，透明球的色散情况与玻璃相仿，问哪种颜色光的 $\alpha$ 角最大，哪种颜色光的 $\alpha$ 角最小？
\end{enumerate}
\onepicture[align=right, w=5cm]{\includesvg{figs/23}}

%% answer: 
\jdanswer{
\begin{enumerate}
	\item
	如图

	\item
	$\alpha = 30^{\circ}$

	\item
	红光 $\alpha$ 最大，紫光 $\alpha$ 最小。

\end{enumerate}
}
%% memo:
\memoanswer{
（1）如图
\onepicture[w=5cm]{\includesvg{figs/23_an}}

（2）由折射定律
$\frac{\sin i}{\sin r} = n$ ①

得 $\sin r = \frac{\sin i}{n} = \frac{\frac{\sqrt{2}}{2}}{\sqrt{2}} = \frac{1}{2}$，$r = 30^{\circ}$
由几何关系及对称性，有

$\frac{\alpha}{2} = r - (i - r) = 2r - i$ ②

$\alpha = 4r - 2i$ ②$'$

以 $r = 30^{\circ}$，$i = 45^{\circ}$ 代入得：$\alpha = 30^{\circ}$

（3）红光 $\alpha$ 最大，紫光 $\alpha$ 最小。

\textbf{评分标准}：（1）3 分：路径 2 分。方向 1 分。（只要反射点在两虚线之间的圆上，路径即算对；只要有一只箭头。方向即算对）

（2）5 分：由折射定律（①式）得 $r = 30^{\circ}$ 得 2 分；由几何关系（②式或②$'$式）得出 $\alpha = 30^{\circ}$ 得 3 分。

（3）2 分：每个结论各 1 分。
}

\item
%% number: 24
%% paperName: 1997年上海高考第 24 题 12 分
%% typeId: 6
%% type: 计算题
%% chapter: 电路
%% point: 等效电路
%% method: 无
%% score: 12
%% degree: 450
%% duplicateId: 0
%% body:
如图所示的电路中，$R_{1} = 3\UO$，$R_{2} = 9\UO$，$R_{3} = 6\UO$，电源电动势 $E = 24\UV$，内阻不计，当电键 $S_{1}$、$S_{2}$ 均开启和均闭合时，灯泡 L 都同样正常发光。
\begin{enumerate}
	\item
	写出两种情况下流经灯泡的电流方向：

	$S_{1}$、$S_{2}$ 均开启时：\tkanswer[2cm]{$b \to a$}；

	$S_{1}$、$S_{2}$ 均闭合时：\tkanswer[2cm]{$a \to b$}。
	\item
	求灯泡正常发光时的电阻 $R$ 和电压 $U$。
\end{enumerate}
\onepicture[align=right, w=5cm]{\includesvg{figs/24}}

%% answer: 
\jdanswer{
\begin{enumerate}
	\item
	$b \to a$，$a \to b$；

	\item
	$R = 3\UO$，$U = 4.8\UV$

\end{enumerate}
}
%% memo:
\memoanswer{
（1）$S_{1}$、$S_{2}$ 均开启时：从 $b \to a$ ①

$S_{1}$、$S_{2}$ 均闭合时：从 $a \to b$ ②

（2）$S_{1}$、$S_{2}$ 均开启时流过灯泡电流

$I_{1} = \frac{E}{R_{1} + R_{2} + R}$ ③

$S_{1}$、$S_{2}$ 均闭合时流过灯泡的电流

$I_{2} = \frac{E}{\frac{R R_{1}}{R + R_{1}} + R_{3}} \cdot \frac{R_{1}}{R + R_{1}}$ ④

同样正常发光，应有

$I_{1} = I_{2}$ ⑤

即

$\frac{E}{R_{1} + R_{2} + R} = \frac{E}{\frac{R R_{1}}{R + R_{1}} + R_{3}} \cdot \frac{R_{1}}{R + R_{1}}$ ⑥

解之，得

$R = \frac{R_{1}(R_{1} + R_{2} - R_{3})}{R_{3}} = \frac{3(3 + 9 - 6)}{6}\UO = 3\UO$

$U = \frac{E}{R_{1} + R_{2} + R}\cdot R = \frac{24}{3 + 9 + 3}\times 3\UV = 4.8\UV$

\textbf{评分标准}：（1）3 分：其中① 1 分② 2 分

（2）9 分：列出③式 1 分，列出④式 2 分。列出⑤式 2 分（直接列出⑥式得 5 分），得出 $R$、$U$ 结果各 2 分。
}

\item
%% number: 25
%% paperName: 1997年上海高考第 25 题 12 分
%% typeId: 6
%% type: 计算题
%% chapter: 气体性质
%% point: 查理定律
%% method: 无
%% score: 12
%% degree: 650
%% duplicateId: 0
%% body:
有人设计了一种测温装置，其结构如图所示。玻璃泡 A 内封有一定量气体，与 A 相连的 B 管插在水银槽中，管内水银面的高度 $x$ 即可反映泡内气体的温度，即环境温度，并可由 B 管上的刻度直接读出。设 B 管的体积与 A 泡的体积相比可略去不计。
\begin{enumerate}
	\item
	在标准大气压下对 B 管进行温度刻度（标准大气压相当于 $76\UcmHg$ 的压强）。已知当温度 $t_{1} = 27\UdC$ 时，管内水银面高度 $x_{1} = 16\Ucm$，此高度即为 $27\UdC$ 的刻度线。问 $t = 0\UdC$ 的刻度线在 $x$ 为多少 $\Ucm$ 处？
	\item
	若大气压已变为相当于 $75\UcmHg$ 的压强，利用该测温装置测量温度时所得读数仍为 $27\UdC$，问此时实际温度为多少？
\end{enumerate}
\onepicture[align=right, h=6cm]{\includesvg{figs/25}}

%% answer: 
\jdanswer{
\begin{enumerate}
	\item
	$x = 21.4\Ucm$

	\item
	$T' = 295\UK = 22\UdC$

\end{enumerate}
}
%% memo:
\memoanswer{
（1）为等容过程，有

$p = \frac{T}{T_{1}}p_{1}$                                                ①

以 $p_{1} = 76 - 16 = 60\UcmHg$，$T_{1} = 273 + 27 = 300\UK$，$T = 273\UK$

代入上式，得 $p = \frac{273}{300}\times 60 = 54.6\UcmHg$

$x = 76 - 54.6 = 21.4\Ucm$

（2）此时 A 泡内气体压强

$p' = 75 - 16 = 59\UcmHg$

实际温度：$T' = \frac{p'}{p}T_{1}$                                                ②

代入数据：$T' = \frac{59}{60}\times 300 = 295\UK = 22\UdC$

\textbf{评分标准}：（1）6 分。知道为等容过程并列出①式得 3 分，得出 $x = 21.4\Ucm$ 得 3 分。

（2）6 分。求出 $p' = 59\UcmHg$ 得 2 分，由公式②求出 $T = 295\UK$（或 $22\UdC$）得 4 分。其中公式②占 2 分，结果占 2 分。
}

\item
%% number: 26
%% paperName: 1997年上海高考第 26 题 14 分
%% typeId: 6
%% type: 计算题
%% chapter: 运动和力
%% point: 动量定理
%% method: 建立模型
%% score: 14
%% degree: 450
%% duplicateId: 0
%% body:
静止在太空中的飞行器上，有一种装置，它利用电场加速带电粒子，形成向外发射的高速粒子流，从而对飞行器产生反冲力，使其获得加速度。已知飞行器质量为 $M$，发射的是 2 价氧离子，发射离子的功率恒为 $P$，加速的电压为 $U$，每个氧离子的质量为 $m$，单位电荷的电量为 $e$。不计发射氧离子后飞行器质量的变化，求：
\begin{enumerate}
	\item
	射出的氧离子速度；
	\item
	每秒钟射出的氧离子数；
	\item
	射出离子后飞行器开始运动时的加速度。
\end{enumerate}

%% answer: 
\jdanswer{
\begin{enumerate}
	\item
	$v = 2\sqrt{\frac{eU}{m}}$

	\item
	$n = \frac{P}{2eU}$

	\item
	$a = \frac{P}{M}\sqrt{\frac{m}{eU}}$

\end{enumerate}
}
%% memo:
\memoanswer{
（1）$qU = \frac{1}{2}mv^{2}$                                                ①

$v = \sqrt{\frac{2qU}{m}} = 2\sqrt{\frac{eU}{m}}$                                                ②

（2）设每秒钟射出的氧离子数为 $n$，有

$P = nqU$                                                ③

得 $n = \frac{P}{qU} = \frac{P}{2eU}$                                                ④

（3）$F = \frac{\Delta p}{\Delta t} = nmv = \frac{P}{2eU}\cdot m\cdot 2\sqrt{\frac{eU}{m}} = P\sqrt{\frac{m}{eU}}$                                                ⑤

$a = \frac{F}{M} = \frac{P}{M}\sqrt{\frac{m}{eU}}$                                                ⑥

\textbf{评分标准}：共 14 分。（1）4 分，写出①得 2 分，得出结果②式 2 分。

（2）5 分。列出③式 3 分，得出④式 2 分。

（3）5 分。得出⑤式 3 分，得出⑥式 2 分。
}

\end{enumerate}

\end{document}
